Στην Άσκηση 7, πραγματοποιήθηκε τυπική επαλήθευση (Formal Verification) χρησιμοποιώντας το εργαλείο Conformal LEC για να επιβεβαιωθεί η λογική ακεραιότητα της σχεδίασης κατά τη διαδικασία της σύνθεσης. Εκτελέστηκαν δύο συγκρίσεις: αρχικά συγκρίθηκε ο κώδικας RTL με το στιγμιότυπο του elaboration (RTL vs Elaborated) και στη συνέχεια ο RTL με το τελικό βελτιστοποιημένο netlist (RTL vs Syn\_Opt). Και στις δύο περιπτώσεις, η διαδικασία ολοκληρώθηκε επιτυχώς με αποτέλεσμα PASS, επιβεβαιώνοντας ότι το κύκλωμα παρέμεινε λειτουργικά αναλλοίωτο.

Συγκεκριμένα, στην τελική σύγκριση (RTL vs Final), η εντολή \texttt{report\_verification} επιβεβαίωσε την ισοδυναμία χωρίς κανένα σφάλμα (0 Non-Equivalent points). Το εργαλείο χρησιμοποίησε ως σημεία σύγκρισης (Compare Points) τις 239 Πρωτεύουσες Εξόδους (Primary Outputs) και τα 2399 Flip-Flops (State Key Points) του σχεδίου, καθώς και 18 σημεία συγχώνευσης (Sequential Merges). Το άθροισμα αυτών $(239+2399+18)$ μας δίνει το συνολικό πλήθος των 2656 Equivalent Points που αναφέρθηκαν στην εντολή \texttt{report\_verification}.